%% kapa: exam=1 round=26 year=2015 period=1 form=A subject=economics subject_ko=경제원론 seq=7 number=47
%% source: Q-Net 감정평가사 기출문제 (cst00309) artlSeq=5138001 file=extracted/1008/5138001/제26회_감정평가사 1차_기출/제26회_감정평가사 1차_기출/1교시_감평_A형.pdf
%% conversion: kordoc 4.24.51 -> tools/gen1.py (auto); review=proofread-sonnet; flags=none
\begin{question}{47}
베이글과 크림치즈는 서로 보완재이고, 베이글과 베이컨은 서로 대체재이다. 베이글의 원료인 밀가루 가격의 급등에 따라 베이글의 생산비용이 상승하였을 때, 각 시장의 변화로 옳지 않은 것은? (단, 베이글, 크림치즈, 베이컨 모두 수요와 공급의 법칙을 따르며, 다른 조건은 일정하다.)
\begin{choices}
  \item 베이글의 가격은 상승한다.
  \item 크림치즈의 거래량은 감소한다.
  \item 크림치즈 시장의 생산자잉여는 감소한다.
  \item 베이컨의 판매수입은 증가한다.
  \item 베이컨 시장의 총잉여는 변함이 없다.
\end{choices}
\end{question}
\answer{5}
